% Generated draft: requires mathematical and source review.
\documentclass[11pt]{article}
\usepackage[margin=25mm]{geometry}
\usepackage{fontspec}
\setmainfont{Latin Modern Roman}
\usepackage{amsmath,amssymb,mathtools,unicode-math}
\setmathfont{Latin Modern Math}
\usepackage{enumitem,needspace,xurl,hyperref}
\setlength{\parindent}{0pt}
\setlength{\parskip}{5pt}
\setlength{\emergencystretch}{3em}
\newcommand{\N}{\mathbb{N}}
\newcommand{\Z}{\mathbb{Z}}
\newcommand{\Q}{\mathbb{Q}}
\newcommand{\R}{\mathbb{R}}
\newcommand{\C}{\mathbb{C}}
\newcommand{\F}{\mathbb{F}}
\newcommand{\A}{\mathbb{A}}
\newcommand{\PP}{\mathbb{P}}
\newcommand{\E}{\mathbb{E}}
\newcommand{\eps}{\varepsilon}
\newcommand{\dd}{\,\mathrm{d}}
\begin{document}
\section*{UnsolvedMath ID: 30006620}
\textbf{Source problem code:} catalogue entry
\subsection{Definitions and mathematical statement}
Let \(C=\{0,1\}^{\N}\) be the Cantor space with its product topology. For a finite binary word \(u\), let \(uC\) denote the cylinder consisting of infinite binary sequences beginning with \(u\). Brin's higher-dimensional Thompson group \(2V\) can be described as the group of homeomorphisms of \(C^2\) given by finite prefix replacements: there are two finite partitions
\[
C^2=\bigsqcup_{i=1}^{m}(u_iC\times v_iC)
=\bigsqcup_{i=1}^{m}(p_iC\times q_iC)
\]
such that the homeomorphism sends
\[
(u_i\alpha,v_i\beta)\longmapsto(p_i\alpha,q_i\beta)
\qquad(\alpha,\beta\in C).
\]
The finite words may be empty, and the group operation is composition.

For a connected, undirected graph \(M\), equip its vertex set with the graph distance \(d_M\), and define
\[
I(x,y)=\{z\in V(M):d_M(x,z)+d_M(z,y)=d_M(x,y)\}.
\]
The graph is a \emph{median graph} if, for every three vertices \(x,y,z\), the intersection
\[
I(x,y)\cap I(y,z)\cap I(z,x)
\]
contains exactly one vertex. Interpret an action on a median graph as an action by graph automorphisms. Such an action has an unbounded orbit if
\[
\exists x\in V(M)\quad \sup_{g\in 2V}d_M(x,gx)=\infty.
\]
For an action on a connected graph, unboundedness of one vertex orbit implies unboundedness of every vertex orbit. No local finiteness, finite-dimensionality, or properness condition is imposed in the supplied question.

A discrete group \(G\) has \emph{Kazhdan's property \((T)\)} if every unitary representation \(\pi:G\to\mathcal U(\mathcal H)\) on a complex Hilbert space that has almost invariant unit vectors also has a nonzero invariant vector. Here almost invariant unit vectors means
\[
\forall F\subseteq G\text{ finite}\;\forall\varepsilon>0\;\exists\xi\in\mathcal H:
\quad \|\xi\|=1,\qquad
\max_{g\in F}\|\pi(g)\xi-\xi\|<\varepsilon,
\]
with the condition on \(F\) understood vacuously when \(F\) is empty. An invariant vector \(\eta\ne0\) satisfies \(\pi(g)\eta=\eta\) for every \(g\in G\).

\textbf{Questions.} For Brin's group \(2V\):
\begin{enumerate}
\item Does there exist a nonempty connected median graph \(M\) and an action \(2V\curvearrowright M\) by graph automorphisms with an unbounded vertex orbit?
\item Does \(2V\) have Kazhdan's property \((T)\)?
\end{enumerate}
These are retained as two questions from the supplied cleaned statement; no equivalence between them is asserted. The supplied literature assessment specifically identifies the property-\((T)\) question as part of Problem 9(i). The raw excerpt begins with an unresolved pronoun and then includes other problems, so the identification of its subject with \(2V\), and the inclusion of the first question, rely on the supplied cleaned statement and curation.
\subsection{Short English statement}
The questions ask whether Brin's group \(2V\) can act on a median graph with an unbounded orbit, and whether every unitary representation of \(2V\) with almost invariant unit vectors must have a nonzero invariant vector. The supplied record labels the problem open; this draft does not independently verify that status.
\subsection{Sources}
\begin{itemize}
\item Supplied local UnsolvedMath registry record: permanent ID \(30006620\), source problem code \texttt{catalogue entry}, title ``Median Actions and Kazhdan Property for Brin Group 2V,'' in the collection ``Oberwolfach Reports Open Problems.'' The record gives the proposed year as 2026 and does not name a proposer.
\item Supplied report citation: \emph{Median Geometry and Applications} (2026), \emph{Oberwolfach Reports}, pp.~487--540, DOI \texttt{10.4171/owr/2026/8}. Supplied source link: \url{https://doi.org/10.4171/owr/2026/8}.
\item Supplied literature source, labelled ``OWR report'': \url{https://ems.press/content/serial-article-files/53603}.
\item Literature assessment date: \texttt{2026-08-22}. This is supplied dated curation, not an independently verified assessment. The curation labels the problem open, identifies the property-\((T)\) question as part of Problem 9(i) for Brin's group \(2V\), and reports that no primary result resolving it was found during that review. The background explicitly calls for independent verification of time-sensitive status and cited claims.
\item The registry supplies the statement-status label \texttt{corrected\_verified} and a verification note asserting that only formula delimiters, notation, grammar, and typography were repaired while preserving mathematical content and scope. These are supplied provenance assertions, not verification performed for this draft.
\end{itemize}
\end{document}
